\documentclass{beamer}

\pdfmapfile{+sansmathaccent.map}


\mode<presentation>
{
  \usetheme{Warsaw} % or try Darmstadt, Madrid, Warsaw, Rochester, CambridgeUS, ...
  \usecolortheme{seahorse} % or try seahorse, beaver, crane, wolverine, ...
  \usefonttheme{serif}  % or try serif, structurebold, ...
  \setbeamertemplate{navigation symbols}{}
  \setbeamertemplate{caption}[numbered]
} 


%%%%%%%%%%%%%%%%%%%%%%%%%%%%
% itemize settings

\definecolor{mypink}{RGB}{255, 30, 80}
\definecolor{mydarkblue}{RGB}{60, 160, 255}
\definecolor{myblue}{RGB}{240, 240, 255}
\definecolor{mygreen}{RGB}{0, 200, 0}
\definecolor{mygreen2}{RGB}{245, 255, 230}
\definecolor{mygray}{gray}{0.8}

\setbeamertemplate{itemize items}[default]

\setbeamertemplate{itemize item}{\color{mygreen}$\blacksquare$}
\setbeamertemplate{itemize subitem}{\color{mydarkblue}$\blacktriangleright$}
\setbeamertemplate{itemize subsubitem}{\color{mygray}$\blacksquare$}



\setbeamercolor{palette quaternary}{fg=white,bg=mydarkblue}
\setbeamercolor{titlelike}{parent=palette quaternary}

\setbeamercolor{palette quaternary2}{fg=black,bg=myblue}
\setbeamercolor{frametitle}{parent=palette quaternary2}



\setbeamerfont{frametitle}{size=\Large,series=\scshape}
\setbeamerfont{framesubtitle}{size=\normalsize,series=\upshape}





%%%%%%%%%%%%%%%%%%%%%%%%%%%%
% block settings

\setbeamercolor{block title}{bg=red!30,fg=black}

\setbeamercolor*{block title example}{bg=mygreen!40!white,fg=black}

\setbeamercolor*{block body example}{fg= black,
bg= mygreen2}


%%%%%%%%%%%%%%%%%%%%%%%%%%%%
% URL settings
\hypersetup{
    colorlinks=false,
    linkcolor=blue,
    filecolor=blue,      
    urlcolor=blue,
}

%%%%%%%%%%%%%%%%%%%%%%%%%%

\renewcommand{\familydefault}{\rmdefault}

\usepackage{amsmath}
\usepackage{mathtools}


\usepackage{subcaption}

\newcommand{\bo}[1] {\mathbf{#1}}
\newcommand{\R} {\mathbb{R}}
\DeclareMathOperator*{\argmin}{arg\,min}


%%%%%%%%%%%%%%%%%%%%%%%%%%%%
% code settings

\usepackage{listings}
\usepackage{color}
% \definecolor{mygreen}{rgb}{0,0.6,0}
% \definecolor{mygray}{rgb}{0.5,0.5,0.5}
\definecolor{mymauve}{rgb}{0.58,0,0.82}
\lstset{ 
  backgroundcolor=\color{white},   % choose the background color; you must add \usepackage{color} or \usepackage{xcolor}; should come as last argument
  basicstyle=\footnotesize,        % the size of the fonts that are used for the code
  breakatwhitespace=false,         % sets if automatic breaks should only happen at whitespace
  breaklines=true,                 % sets automatic line breaking
  captionpos=b,                    % sets the caption-position to bottom
  commentstyle=\color{mygreen},    % comment style
  deletekeywords={...},            % if you want to delete keywords from the given language
  escapeinside={\%*}{*)},          % if you want to add LaTeX within your code
  extendedchars=true,              % lets you use non-ASCII characters; for 8-bits encodings only, does not work with UTF-8
  firstnumber=0000,                % start line enumeration with line 0000
  frame=single,	                   % adds a frame around the code
  keepspaces=true,                 % keeps spaces in text, useful for keeping indentation of code (possibly needs columns=flexible)
  keywordstyle=\color{blue},       % keyword style
  language=Octave,                 % the language of the code
  morekeywords={*,...},            % if you want to add more keywords to the set
  numbers=left,                    % where to put the line-numbers; possible values are (none, left, right)
  numbersep=5pt,                   % how far the line-numbers are from the code
  numberstyle=\tiny\color{mygray}, % the style that is used for the line-numbers
  rulecolor=\color{black},         % if not set, the frame-color may be changed on line-breaks within not-black text (e.g. comments (green here))
  showspaces=false,                % show spaces everywhere adding particular underscores; it overrides 'showstringspaces'
  showstringspaces=false,          % underline spaces within strings only
  showtabs=false,                  % show tabs within strings adding particular underscores
  stepnumber=2,                    % the step between two line-numbers. If it's 1, each line will be numbered
  stringstyle=\color{mymauve},     % string literal style
  tabsize=2,	                   % sets default tabsize to 2 spaces
  title=\lstname                   % show the filename of files included with \lstinputlisting; also try caption instead of title
}

%%%%%%%%%%%%%%%%%%%%%%%%%%%%
% tikz settings

\usepackage{tikz}
\tikzset{every picture/.style={line width=0.75pt}}

%%%%%%%%%%%%%%%%%%%%%%%%%%%%

\usepackage{qrcode}



\title{Observers}
\subtitle{Control Theory, Lecture 10}
\author{by Sergei Savin}
\centering
\date{Spring 2021}



\begin{document}
\maketitle



\begin{frame}{Content}
\begin{itemize}
\item Measurement
\begin{itemize}
    \item How do we know the state?
    \item Why information is imperfect?
    \item Definition
\end{itemize}
\item Observation
\begin{itemize}
    \item Using the knowledge about dynamics
    \item Observer
    \item Observer gains
\end{itemize}
\item Observer Design
\item Observation and Control
\begin{itemize}
    \item LTI
    \item Stability analysis
    \item Change of variables
    \item Upper triangular form
    \item Separation principle
\end{itemize}
\end{itemize}
\end{frame}




\begin{frame}{Measurement}
\framesubtitle{How do we know the state?}
\begin{flushleft}

Before we considered systems and control laws of the following type:

\begin{equation}
\begin{cases}
\dot {\bo{x}} = \bo{A} \bo{x} + \bo{B} \bo{u}\\
\bo{u} = \bo{K} \bo{x}
\end{cases}
\end{equation}

But when we implement that control law, how do we know the current value of $\bo{x}$? Previously we always took it from simulation. 

\bigskip

In practice, we take it from \emph{measurement}.

\end{flushleft}
\end{frame}

\begin{frame}{Measurement}
\framesubtitle{Why information is imperfect?}
\begin{flushleft}

There are a number of reasons why we can not directly measure the state of the system. Here are some:

\begin{itemize}
\item Digital measurements are done in discrete time intervals;
\item Unpredicted events (faults, collisions, etc.);
\item Un-modelled kinematics or dynamics (links bending, gear box backlash,  friction, etc.) making the very definition of the state disconnected from reality;
\item Lack of sensors;
\item Imprecise, nonlinear and biased sensors;
\item Physics, quantum-scale effects and alike;
\end{itemize}

\end{flushleft}
\end{frame}

\begin{frame}{Measurement}
\framesubtitle{Definition}
\begin{flushleft}

Let us introduce new notation. Assume we have an LTI system of the following form:

\begin{equation}
\begin{cases}
\dot {\bo{x}} = \bo{A} \bo{x} + \bo{B} \bo{u} \\
\bo{y} = \bo{C} \bo{x} \\
\bo{u} = \bo{g}(\hat{\bo{x}}, t)\\
\hat{\bo{x}} = \bo{o}(\cdot)
\end{cases}
\end{equation}

Then:

\begin{itemize}
\item $\bo{x}$ is the state (actual or true state)
\item $\bo{y}$ is the output (actual or true output)
% \item $\Tilde{\bo{y}}$ is the measured output
\item $\hat{\bo{x}}$ is the estimated (observed) state
\item $\hat{\bo{y}}$ is the estimated (observed) output
\end{itemize}

Notice that we never know true state $\bo{x}$, and therefore for the control purposes we have to use the estimated state $\hat{\bo{x}}$.

\end{flushleft}
\end{frame}



\begin{frame}{Observation}
\framesubtitle{Using the knowledge about dynamics}
\begin{flushleft}

Let us consider autonomous dynamical system
\begin{equation}
\label{eq:LTI}
\begin{cases}
\dot {\bo{x}} = \bo{A} \bo{x} + \bo{B} \bo{u} \\
\bo{y} = \bo{C} \bo{x}
\end{cases}
\end{equation}
%
with measurements $\bo{y}$. We want to get as good an estimate of the state $\hat{\bo{x}}$ as we can.

\bigskip

First note: dynamics should also hold for our observed state:
\begin{equation}
\hat{\dot {\bo{x}}} = \bo{A} \hat{\bo{x}} + \bo{B} \bo{u}
\end{equation}
%
Therefore if we know the initial conditions of our system exactly, and we know our model exactly, we can find exact state of the system without using measurement $\bo{y}$. We can call it an open loop observation. Unfortunately, we know neither the model nor the initial conditions precisely.


\end{flushleft}
\end{frame}





\begin{frame}{Observation}
\framesubtitle{Observer}
\begin{flushleft}

We propose \emph{observer} that takes into account measurements in a linear way:

\begin{equation}
\label{eq:Observer}
\hat{\dot {\bo{x}}} = \bo{A} \hat{\bo{x}} + \bo{B} \mathbf u + \bo{L}(\mathbf y - \bo{C} \hat{\bo{x}})
\end{equation}
%
with measurements $\bo{y}$. With this observer, we want to get as good estimate of the state $\hat{\bo{x}}$ as we can.

\bigskip

Let's define state estimation error as $\varepsilon = \hat{\bo{x}} - \bo{x}$. We can subtract \eqref{eq:LTI} from \eqref{eq:Observer}, to get \emph{observer error dynamics}:

\begin{equation}
\hat{\dot {\bo{x}}} - \dot {\bo{x}}= 
\bo{A} \hat{\bo{x}} - \bo{A} \bo{x} + 
\bo{L}(\mathbf y - \bo{C} \hat{\bo{x}})
\end{equation}

\begin{equation}
\dot {\varepsilon}= 
(\bo{A} - \bo{L} \bo{C}) \varepsilon
\end{equation}

\end{flushleft}
\end{frame}





\begin{frame}{Observation}
\framesubtitle{Observer gains}
\begin{flushleft}

The observer $\dot {\varepsilon}= 
(\bo{A} - \bo{L} \bo{C}) \varepsilon$ is \emph{stable} (i.e., the state estimation error tends to zero), as long as the following matrix has eigenvalues with negative real parts:

\[
\bo{A} - 
\bo{L} \bo{C} < 0
\]

We need to find $\bo{L}$. Let us observe the key difference between observer design and controller design:

\bigskip

\begin{itemize}
    \item Controller design: find such $\bo{K}$ that $\bo{A} - \bo{B} \bo{K} < 0$.
    \item Observer design: find such $\bo{L}$ that: $\bo{A} - \bo{L} \bo{C} < 0$
\end{itemize}

\bigskip

We have instruments for finding $\bo{K}$, what about $\bo{L}$?

\end{flushleft}
\end{frame}


\begin{frame}{Observer Design}
\framesubtitle{General case: design via Riccati eq.}
\begin{flushleft}

In general, we can observe that if $\bo{A} - \bo{L} \bo{C}$ is negative-definite, then $(\bo{A} - 
\bo{L} \bo{C})^{\top}$ is negative-definite too (by definition of the negative-definiteness). 

\bigskip

Therefore, we can solve the following \emph{dual problem}:

\begin{itemize}
    \item find such $\bo{L}$ that $\bo{A}^{\top} - 
\bo{C}^{\top} \bo{L}^{\top} < 0$.
\end{itemize}

\bigskip

The dual problem is \emph{equivalent} to the control design problem. We can solve it by producing and solving algebraic Riccati equation, as in the LQR formulation. In pseudo-code it can be represented the following way:

\bigskip

$\bo{L}^{\top}$ \texttt{= lqr}($\bo{A}^{\top}$, $\bo{C}^{\top}$, $\mathbf Q$, $\mathbf R$).

where $\mathbf Q$ and $\mathbf R$ are weight  matrices, determining the "sensitivity" or "aggressiveness" of the observer.


\end{flushleft}
\end{frame}




\begin{frame}{Observation and Control}
\framesubtitle{LTI}
\begin{flushleft}

Thus we get dynamics+observer combination:

\begin{equation}
\begin{cases}
\dot {\bo{x}} = \bo{A} \bo{x} + \bo{B} \bo{u} \\
\hat{\dot {\bo{x}}} = \bo{A} \hat{\bo{x}} + \bo{B} \mathbf u + \bo{L}(\mathbf y - \bo{C} \hat{\bo{x}})\\
\bo{y} = \bo{C} \bo{x} \\
\bo{u} = -\bo{K} (\hat{\bo{x}} - \bo{x}^*(t)) + \bo{u}^*(t)
\end{cases}
\end{equation}

\bigskip

where $\bo{A} - \bo{B} \bo{K} < 0$ and $\bo{A}^{\top} - 
\bo{C}^{\top} \bo{L}^{\top} < 0$.


\end{flushleft}
\end{frame}




\begin{frame}{Observation and Control}
\framesubtitle{Stability analysis}
\begin{flushleft}

When we want to stabilize the origin, the dynamics+observer combination can be represented as a single LTI:

\begin{equation}
\begin{cases}
\dot {\bo{x}} = \bo{A} \bo{x} + \bo{B} \bo{u} \\
\hat{\dot {\bo{x}}} = \bo{A} \hat{\bo{x}} + \bo{B} \mathbf u + \bo{L}(\mathbf y - \bo{C} \hat{\bo{x}})\\
\bo{y} = \bo{C} \bo{x} \\
\bo{u} = -\bo{K} \hat{\bo{x}}
\end{cases}
\end{equation}

In matrix form it becomes:

\begin{equation}
\begin{bmatrix}
\dot {\bo{x}} \\
\hat{\dot {\bo{x}}}
\end{bmatrix}
=
\begin{bmatrix}
\bo{A} & -\bo{B}\bo{K} \\
\bo{L}\bo{C} & (\bo{A} - \bo{B}\bo{K}-\bo{L}\bo{C})
\end{bmatrix}
\begin{bmatrix}
\bo{x} \\
\hat{\bo{x}}
\end{bmatrix}
\end{equation}

\bigskip

We can't directly reason about eigenvalues of this matrix. Next slide will show a way to do it with a change of variables.

\end{flushleft}
\end{frame}




\begin{frame}{Observation and Control}
\framesubtitle{Change of variables}
\begin{flushleft}

Let us use the following substitution: $\bo{e} = \bo{x} - \hat{\bo{x}}$, which implies $\hat{\bo{x}} = \bo{x} - \bo{e}$:

Our system had form:

\begin{equation}
\begin{cases}
\dot {\bo{x}} = \bo{A} \bo{x} - \bo{B}\bo{K} \hat{\bo{x}} \\
\hat{\dot {\bo{x}}} = \bo{A} \hat{\bo{x}} - \bo{B}\bo{K} \hat{\bo{x}} + \bo{L}(\bo{C} \bo{x} - \bo{C} \hat{\bo{x}})
\end{cases}
\end{equation}

Since $\dot{\bo{e}} = \dot{\bo{x}} - \hat{\dot{\bo{x}}}$, we get:
%
\[
\dot{\bo{e}} = 
\bo{A} \bo{x} - \bo{B}\bo{K} \hat{\bo{x}} - 
(\bo{A} \hat{\bo{x}} - \bo{B}\bo{K} \hat{\bo{x}} + \bo{L}(\bo{C} \bo{x} - \bo{C} \hat{\bo{x}}))
\]
%
\[
\dot{\bo{e}} = 
\bo{A} (\bo{x} - \hat{\bo{x}})  - \bo{L}\bo{C}(\bo{x} - \hat{\bo{x}})
\]
%
\[
\dot{\bo{e}} = 
(\bo{A}  - \bo{L}\bo{C})\bo{e}
\]

Equation for $\dot {\bo{x}}$ takes form:

\[
\dot {\bo{x}} = (\bo{A}-\bo{B}\bo{K}) \bo{x} +  \bo{B}\bo{K}\bo{e}
\]


\end{flushleft}
\end{frame}




\begin{frame}{Observation and Control}
\framesubtitle{Upper triangular form}
\begin{flushleft}

Collecting $\dot {\bo{x}}$ and $\dot{\bo{e}}$ we get:

\begin{equation}
\begin{cases}
\dot {\bo{x}} = (\bo{A}-\bo{B}\bo{K}) \bo{x} +  \bo{B}\bo{K}\bo{e} \\
\dot{\bo{e}} = 
(\bo{A}  - \bo{L}\bo{C})\bo{e}
\end{cases}
\end{equation}

In matrix form it becomes:

\begin{equation}
\begin{bmatrix}
\dot {\bo{x}} \\
\dot{\bo{e}}
\end{bmatrix}
=
\begin{bmatrix}
(\bo{A}-\bo{B}\bo{K}) & \bo{B}\bo{K} \\
0 & (\bo{A}  - \bo{L}\bo{C})
\end{bmatrix}
\begin{bmatrix}
\bo{x} \\
\bo{e}
\end{bmatrix}
\end{equation}

Eigenvalues of a upper block-triangular matrices equal to the union of the eigenvalues of the blocks on the main diagonal. Hence here, the eigenvalues of the system are equal to the union of eigenvalues of $(\bo{A}-\bo{B}\bo{K})$ and $(\bo{A}  - \bo{L}\bo{C})$. 

\end{flushleft}
\end{frame}



\begin{frame}{Observation and Control}
\framesubtitle{Separation principle}
\begin{flushleft}
 
Since the eigenvalues of the system are equal to the union of eigenvalues of $(\bo{A}-\bo{B}\bo{K})$ and $(\bo{A}  - \bo{L}\bo{C})$, we can make the following observation:

\bigskip

\begin{alertblock}{Separation principle}
As long as the observer and the controller are stable independently, the overall system is stable too. This is called \emph{separation principle}.
\end{alertblock}

\end{flushleft}
\end{frame}




% \begin{frame}{Observation and Control}
% \framesubtitle{Affine case}
% \begin{flushleft}


% Affine case is almost the same:

% \begin{equation}
% \begin{cases}
% \dot {\bo{x}} = \bo{A} \bo{x} + \bo{B} \bo{u} + \bo{c}\\
% \hat{\dot {\bo{x}}} = \bo{A} \hat{\bo{x}} + \bo{B} \mathbf u + \bo{L}(\mathbf y - \bo{C} \hat{\bo{x}})  + \bo{c} \\
% \bo{y} = \bo{C} \bo{x} \\
% \bo{u} = -\bo{K} (\hat{\bo{x}} - \bo{x}^*(t)) + \bo{u}^*(t)
% \end{cases}
% \end{equation}

% \bigskip

% where $\bo{A} - \bo{B} \bo{K} < 0$ and $\bo{A}^{\top} - 
% \bo{c}^{\top} \bo{L}^{\top} < 0$.


% \end{flushleft}
% \end{frame}





\begin{frame}{Thank you!}
\centerline{Lecture slides are available via Moodle.}
\bigskip
\centerline{You can help improve these slides at:}
\centerline{\href{https://github.com/SergeiSa/Control-Theory-Slides-Spring-2021}{github.com/SergeiSa/Control-Theory-Slides-Spring-2021}}

\includegraphics[width=1.6in]{qrcode.png}
\bigskip

\centerline{Check Moodle for additional links, videos, textbook suggestions.}

\end{frame}

\end{document}
